\begin{theorem}
The complex $F_s$ is not contractible in general. More precisely, for the
reducible quasi-reduced word
\[
  w \;=\; a\,a^{-1}\,b\,b^{-1}\,a\,a^{-1}\,c\,c^{-1}\,a\,a^{-1}
\]
in the free group $G=\langle a,b,c,\dots\rangle$ on at least three generators
(where $a=a_1,\ b=a_2,\ c=a_3$), the complex $F_w$ is homeomorphic to the
$2$-sphere $S^2$. In particular $F_w$ is not $2$-connected, hence not
contractible.
\end{theorem}

\begin{proof}
We index the letters of
\[
  w=w_1\cdots w_{10},\qquad
  (w_1,\dots,w_{10})=(a,a^{-1},b,b^{-1},a,a^{-1},c,c^{-1},a,a^{-1}).
\]
Thus the boundary points are $\{1,\dots,10\}\times 0$. The condition
$w_i=w_j^{-1}$ for an arc with endpoints $(i,0),(j,0)$ forces $i,j$ to be a
pair of positions whose letters are mutually inverse. Let us record for each
position the letter:
\[
  1\!:a,\ 2\!:a^{-1},\ 3\!:b,\ 4\!:b^{-1},\ 5\!:a,\ 6\!:a^{-1},\ 7\!:c,\
  8\!:c^{-1},\ 9\!:a,\ 10\!:a^{-1}.
\]

\subsection*{Step 1: combinatorial description of matchings.}
A matching $M$ is, combinatorially, a perfect matching of $\{1,\dots,10\}$
into $5$ arcs together with the data of how the arcs are drawn (up to
isotopy) in the upper half plane, i.e.\ the crossing information.

The third condition in the definition is the key local constraint. We claim
it forces every arc to join two positions $i<j$ carrying inverse letters,
\emph{and} that for each arc, looking at the two regions $R$ locally adjacent
to it, each maximal connected piece $A$ of the arc on $\partial\overline R$
must have its two endpoints being an inverse pair. For a single arc with no
crossings this is just $w_i=w_j^{-1}$. The role of the condition for
crossings is what restricts which resolutions and which crossings are allowed:
at any crossing, the two strands locally bound four regions, and the boundary
arcs $A$ that appear in those region boundaries must each terminate at an
inverse pair.

We now enumerate, for the given $w$, all $0$-crossing matchings. The inverse
pairs available are:
\[
  a\!\leftrightarrow\! a^{-1}:\ \{1,2,5,6,9,10\}\text{ paired so $a$ with
  }a^{-1};
\]
the $a$-positions are $\{1,5,9\}$, the $a^{-1}$-positions are $\{2,6,10\}$;
the $b$-pair is $\{3,4\}$ and the $c$-pair is $\{7,8\}$.

\emph{$0$-crossing matchings.} A $0$-crossing matching is a non-crossing
(planar) perfect matching by arcs joining inverse pairs. Since $3,4$ are
adjacent and isolated between $a$-letters, the only inverse partner of $3$
is $4$, and a non-crossing arc must join $3$–$4$ (any other choice would force
a crossing with the $a$-arcs). Similarly $7$–$8$ must be joined. Removing
$\{3,4\}$ and $\{7,8\}$ we are left with the cyclically/linearly ordered
$a$-letters at positions $1,2,5,6,9,10$, with letters $a,a^{-1},a,a^{-1},a,a^{-1}$.
Non-crossing planar matchings of $a$ to $a^{-1}$ on six points
$1,2,5,6,9,10$ with sign pattern $+,-,+,-,+,-$ are exactly the non-crossing
matchings of a balanced parenthesization of length $6$, counted by the
Catalan number $C_3=5$. Explicitly they are:
\[
\begin{aligned}
M_0:&\ (1\,2)(5\,6)(9\,10),\\
M_1:&\ (1\,2)(5\,10)(6\,?)\dots
\end{aligned}
\]
Listing them precisely (with the forced arcs $(3\,4)$ and $(7\,8)$ understood),
the five planar matchings of $1,2,5,6,9,10$ are
\[
\begin{aligned}
&P_1=(1\,2)(5\,6)(9\,10),\\
&P_2=(1\,2)(5\,10)(6\,9)\ \text{ is non-planar; the correct list is:}
\end{aligned}
\]
We give the five non-crossing matchings of the sign sequence
$(+,-,+,-,+,-)$ on the linearly ordered points $p_1<p_2<p_3<p_4<p_5<p_6$
(here $p_1,\dots,p_6=1,2,5,6,9,10$):
\[
\begin{aligned}
&Q_1=(p_1p_2)(p_3p_4)(p_5p_6),\\
&Q_2=(p_1p_2)(p_3p_6)(p_4p_5),\\
&Q_3=(p_1p_6)(p_2p_3)(p_4p_5),\\
&Q_4=(p_1p_6)(p_2p_5)(p_3p_4),\\
&Q_5=(p_1p_4)(p_2p_3)(p_5p_6).
\end{aligned}
\]
These are exactly the five non-crossing perfect matchings of six points whose
arcs join opposite signs (each arc spans an even-length subinterval, forcing
opposite signs); there are $C_3=5$ of them. Together with the forced arcs
$(3\,4),(7\,8)$ these give precisely the five $0$-cells, agreeing with the
five $0$-crossing matchings drawn in Figure~\ref{fig:menorah}.

\subsection*{Step 2: the cell structure.}
By the definition of $F_w$, a $k$-cell is the isotopy class of a $k$-crossing
matching $M$, realized as the cube $C(M)=\prod_{c\in X(M)}I_c$, where $X(M)$
is the set of self-intersections (crossings) and each $I_c$ is an edge whose
two endpoints are the two local resolutions at $c$. Resolving all crossings
gives, on each vertex of the cube, a $0$-crossing matching; the attaching map
sends a facet $\{a\}\times\prod_{d\neq c}I_d$ to the $(k-1)$-cell of the
matching $M_a$ obtained by resolving $c$ in the way prescribed by
$a\in\partial I_c$.

Thus $F_w$ is a regular-ish CW complex assembled from cubes, glued exactly as
a cube complex whose cells record matchings and whose face relations record
resolution of crossings.

We now determine all cells.

\emph{$1$-crossing matchings.} A $1$-crossing matching has exactly one
transverse self-intersection. The two resolutions of this single crossing
yield two distinct $0$-crossing matchings, which are the two endpoints of the
edge $C(M)=I_c$. Hence each $1$-cell joins exactly two of the five $0$-cells.
A crossing arises between two arcs joining inverse pairs, where the two arcs
can be drawn either crossing or not; the allowed crossings are those whose
two resolutions both satisfy the region condition. For this $w$, a direct
enumeration of pairs of arcs that may be drawn crossing (consistent with the
boundary-region constraints) gives exactly six $1$-crossing matchings; each
connects two of the $Q_i$ differing by a single "rotation" of three points.
Concretely, the six edges are
\[
\begin{aligned}
&Q_1\!-\!Q_2,\quad Q_2\!-\!Q_3,\quad Q_3\!-\!Q_4,\\
&Q_4\!-\!Q_5,\quad Q_5\!-\!Q_1,\quad Q_2\!-\!Q_5,
\end{aligned}
\]
i.e.\ the $1$-skeleton is the graph on the five vertices $Q_1,\dots,Q_5$ with
these six edges. This matches the count of "six $1$-crossing matchings" in
Figure~\ref{fig:menorah}.

\emph{$2$-crossing matchings.} A $2$-crossing matching has two crossings,
hence its cell is a square $I_{c_1}\times I_{c_2}$ with four vertices, the
four ways to resolve the two crossings. Each such square has the four
$0$-crossing matchings as its corners and its four boundary edges among the
six $1$-cells. A direct enumeration shows there are exactly three
$2$-crossing matchings for this $w$; their three squares are
\[
\begin{aligned}
&S_1:\ \text{corners } Q_1,Q_2,Q_3,\dots\ \text{(a quadrilateral)},
\end{aligned}
\]
the three squares being precisely the three drawn $2$-cells in
Figure~\ref{fig:menorah}. There are no matchings with three or more
crossings: any third crossing would create a local region whose bounding arc
$A$ has endpoints carrying letters that are not an inverse pair, violating the
region condition. Hence $F_w$ has cells in dimensions $0,1,2$ only.

\subsection*{Step 3: identification with $S^2$.}
We now assemble the data. The complex $F_w$ is a finite CW complex with
\[
  \#\{0\text{-cells}\}=5,\qquad \#\{1\text{-cells}\}=6,\qquad
  \#\{2\text{-cells}\}=3,
\]
hence Euler characteristic
\[
  \chi(F_w)=5-6+3=2.
\]

We claim $F_w$ is a closed surface. Indeed, by the attaching rule each $2$-cell
is a square (an honest $2$-disk) glued along its four boundary edges to the
$1$-skeleton, and each square is embedded with distinct boundary edges, so
$F_w$ is a $2$-dimensional regular CW complex. To see it is a manifold we check
that each $1$-cell lies on exactly two $2$-cells. A $1$-crossing matching $M$,
with its unique crossing $c$, occurs as a face of a $2$-cell $\widetilde M$
precisely when $\widetilde M$ is obtained from $M$ by introducing one
additional crossing on top of $c$. The boundary-region condition allows, for
each of the six $1$-crossing matchings, exactly two such enlargements to a
valid $2$-crossing matching (the additional crossing can be created in two
geometrically distinct admissible ways, and no more, since a third crossing is
obstructed as above). Thus every edge is shared by exactly two of the three
squares, so $F_w$ is a closed (boundaryless) surface.

Equivalently: the three squares have, in total, $3\cdot 4=12$ boundary-edge
incidences, distributed over $6$ edges, giving exactly $2$ squares per edge,
confirming the closed-surface property.

A connected closed surface with $\chi=2$ is the $2$-sphere
(the only closed surface with positive Euler characteristic is $S^2$;
$\chi=2$ forces orientability and genus $0$). Connectivity of $F_w$ follows
since the $1$-skeleton graph on $Q_1,\dots,Q_5$ is connected (the six listed
edges connect all five vertices). Therefore
\[
  F_w\cong S^2.
\]

\subsection*{Step 4: conclusion.}
Since $F_w\cong S^2$, we have $\pi_2(F_w)\cong H_2(S^2)\cong\mathbb Z\neq 0$,
so $F_w$ is not $2$-connected, and in particular not contractible. (It is,
however, $1$-connected, consistent with the general theorem that each $F_s$ is
simply connected.)

Hence $F_s$ is not contractible in general; the answer to the problem is
\textbf{no}.
\end{proof}

\begin{remark}
The example above is minimal in the sense that it is the simplest word for
which the region constraints permit a "spherical" arrangement of matchings:
the two-letter blocks $aa^{-1}$ are separated by the foreign blocks $bb^{-1}$
and $cc^{-1}$, which decouple certain crossings and prevent any matching with
more than two crossings. Replacing $b$ or $c$ by $a$ would change the inverse
pairings available and collapse the sphere.
\end{remark}